\documentclass{AuthorGuide/ijitcs}

\usepackage[T1]{fontenc}
\usepackage[utf8]{inputenc}
\usepackage{amsmath,amssymb,array,booktabs,graphicx,tabularx}
\usepackage[hidelinks]{hyperref}

\graphicspath{{AuthorGuide/}{../results/plots/}}

\newcommand{\fittableinput}[1]{%
\begingroup
\setlength{\tabcolsep}{4pt}%
\resizebox{\textwidth}{!}{\input{#1}}%
\endgroup
}

\IjitcsSet{
    title = {Reliability of Tabular Machine Learning in Operational Data Pipelines: Structured Corruption Stress Testing for Decision Support and Expert Systems},
    volume = {XX},
    issue = {XX},
    year = {2026},
    month = {Month},
    startpage = {1},
    endpage = {19},
    doi = {TBD},
    abstract = {Modern decision support systems increasingly rely on data-driven pipelines that integrate machine learning models operating on large-scale tabular data. In operational environments, these pipelines are vulnerable to data quality problems such as missing values, categorical drift, numeric inconsistencies, and other structured corruptions that can degrade downstream decision quality. This paper presents a reproducible stress-testing framework for evaluating the reliability of tabular machine-learning pipelines under controlled corruption scenarios. The benchmark studies twelve binary OpenML datasets, three common model families, two categorical-encoding policies, multiple calibration methods, and threshold-governance strategies. The design also emphasizes traceable artifacts and repeatable evaluation settings for peer review. Results show that structured missingness is the most harmful single corruption family in this benchmark, while compound missingness plus numerical inconsistency produces the largest average degradation in discrimination and thresholded performance. Calibration improves probability quality under stress, and adaptive threshold policies reduce expected decision cost relative to fixed-threshold deployment. The study provides an auditable methodology and practical design guidance for building more resilient decision support and expert systems.},
    keywords = {data pipelines, system reliability, decision support systems, data quality, robust machine learning, data corruption},
    received = {Submitted March 28, 2026; Revised: N/A; Accepted: N/A; Published: N/A}
}

\author*{Kyriakos Kartas}{Independent Researcher, Greece}{mail@kkartas.gr}{https://orcid.org/0009-0001-6477-4676}

\thispagestyle{thefirstpage}
\begin{document}

\maketitle

\section{Introduction}
\subsection{Problem Statement}
Operational machine learning deployments often assume that inference-time data follow the same process as model development data. In tabular systems, this assumption is frequently violated by upstream pipeline defects, including join amplification, feature missingness, categorical domain drift, and numeric transformation inconsistencies. These structured operational corruptions are high impact in decision support and expert systems because probabilistic predictions are mapped to threshold-triggered actions.

Many production teams still report model quality primarily through clean-test discrimination metrics. That practice is useful for model selection, but it can be operationally incomplete. A classifier can preserve acceptable ranking behavior under certain data degradations while simultaneously losing probability calibration and policy utility. In systems where a probability score triggers business, clinical, or operational actions, this mismatch is not a marginal issue. A small calibration error can shift the decision boundary enough to create a materially different false-positive/false-negative profile under fixed thresholding.

The central premise of this manuscript is therefore practical: reliability is an end-to-end property of the model-plus-pipeline system. A model with stable clean-test AUC may still be brittle in deployment when data engineering assumptions are violated. Conversely, some models may show moderate discrimination drift but still remain controllable if calibration and threshold governance are designed explicitly. This paper addresses that gap with a benchmark that integrates data corruption design, metric-level reliability analysis, and policy-level utility assessment.

\subsection{Study Context}
The operational framing is intentionally grounded in common tabular deployment architectures. In most organizations, tabular prediction systems consume features created by extraction and transformation logic across heterogeneous data stores. Feature pipelines may involve slowly changing dimensions, delayed updates, schema migrations, category recoding, vendor integrations, and floating-point unit conversions. Each stage can introduce structured, non-i.i.d. perturbations that are not well represented by generic synthetic noise. By modeling these perturbations as explicit corruption families, this work prioritizes failure modes that practitioners encounter in real systems audits.

A second motivation is reproducibility at artifact level rather than narrative level. The benchmark used in this manuscript is not presented as a one-off analysis script; it is represented by deterministic seeds, explicit corruption severity grids, machine-readable summary tables, and generated figure artifacts. This matters for two reasons. First, it allows independent reruns with equivalent protocol semantics. Second, it supports paper-to-code traceability, where a claim in the manuscript can be traced to a specific CSV table and, when needed, to the long-form raw artifact.

This paper emphasizes applied rigor: transparent design choices, explicit uncertainty reporting, and deployment-oriented interpretation. Instead of claiming universal robustness for a model class, we evaluate where and how reliability degrades, which interventions remain effective under stress, and which governance controls are likely to provide the strongest risk reduction in operational settings.

The benchmark scope reflects this applied emphasis. We evaluate 12 binary OpenML datasets, three model families that are widely used in tabular production systems, two categorical encoding policies, and multiple random seeds. The corruption protocol includes four single families (C1--C4) and one compound family (C2+C4), with severity sweeps and severe-versus-baseline contrasts. Outcomes are reported jointly on discrimination (AUC and F1), calibration quality (ECE and calibration diagnostics), and policy-level expected decision cost under explicit false-positive/false-negative asymmetry.

The intended contribution is not a new learning algorithm. Instead, the contribution is a complete reliability-evaluation procedure that can be reused as a deployment readiness audit. The procedure links structured operational corruptions to measurable reliability changes, interprets those changes with multiplicity-aware statistics, and translates metric shifts into policy consequences through threshold governance analysis. This linkage is often missing in conventional benchmark reporting, where calibration and operating-point policy are treated as secondary or omitted entirely.

From a governance perspective, the manuscript advances a three-layer reliability view. Layer one is data quality governance (monitoring completeness, domain integrity, and numeric consistency). Layer two is probability governance (calibration quality and reliability curves). Layer three is policy governance (threshold design aligned with cost asymmetry). Each layer can partially compensate for weaknesses in another, but none can safely replace the others. The empirical sections demonstrate this interaction rather than assuming it.

The practical implication for decision support and expert systems is straightforward. Reliability claims should not be based on clean discrimination alone. A deployment recommendation is credible only when the system is evaluated under realistic corruption stress and when mitigation controls are tested at both probability and policy levels.

\subsection{Research Objective}
The study addresses five questions:
\begin{enumerate}
\item Which corruption families induce the largest reliability degradation?
\item Are compound corruptions stronger than single-family perturbations?
\item How heterogeneous are effects across datasets, models, and encoding policies?
\item Which calibration methods remain stable under severe corruption?
\item How much risk reduction can be obtained through threshold governance?
\end{enumerate}

\subsection{Contribution Statement}
\begin{itemize}
\item A reproducible stress benchmark for tabular reliability on 12 OpenML datasets \cite{vanschoren2014openml}.
\item Structured operational corruptions C1--C4 and compound C2+C4 with severity control and realized diagnostics.
\item Joint evaluation of discrimination (AUC, F1), calibration (ECE), and policy-level expected cost.
\item Predefined setting-level and dataset-level paired inference with Benjamini--Hochberg multiplicity control \cite{benjamini1995fdr}, reported with explicit dependence and power-scope caveats.
\item Deployment-oriented guidance: integrated governance of data quality, calibration, and decision thresholds.
\end{itemize}

\loadmaingeometry

\section{Related Work}
Operational ML failures are commonly driven by data and pipeline issues rather than only algorithmic misspecification \cite{sculley2015hidden}. Our corruption design focuses on this operational perspective through structured perturbations that emulate realistic tabular pipeline faults.
More generally, reliability under non-stationarity has long been studied through dataset-shift and concept-drift lenses \cite{quionero2009datasetshift,moreno2012datasetshift,gama2014conceptdrift}. We use that framing at deployment interface level: corruption families are treated as explicit operational shift channels (missingness, category novelty, numerical inconsistency, and cohort reweighting) rather than as generic i.i.d. noise.

\subsection{Operational Data Reliability and Technical Debt}
The literature on ML system technical debt emphasizes that model quality cannot be isolated from pipeline assumptions, data dependencies, and serving-time infrastructure behavior \cite{sculley2015hidden}. In practice, model artifacts are embedded in larger systems with feature stores, ETL dependencies, data contracts, and service orchestration. Reliability degradations in these settings are often not adversarial in the security sense; instead, they arise from process drift, schema inconsistency, and integration errors that accumulate gradually.

Tabular pipelines are particularly sensitive to this class of faults. Features are often assembled from heterogeneous upstream sources with varied update cadences and quality controls. Joins can over-amplify specific cohorts when key cardinalities are violated. Missingness can be introduced by late-arriving tables, upstream null-casting, or domain-specific masking rules. Categorical drift can emerge when source systems emit new vendor codes, legacy encodings, or transformed identifiers. Numerical inconsistency can result from unit changes, clipping, rounding, and scale shifts induced by cross-system conversions.

Although these issues are broadly recognized by practitioners, many empirical ML studies still evaluate robustness with perturbations that are either unstructured or detached from observed data engineering failure modes. This manuscript addresses that gap by centering structured operational corruptions as first-class benchmark objects. The design objective is not to reproduce any one production incident exactly, but to create stress families that preserve operational plausibility and can be parameterized with controlled severity.

An additional concern in operational studies is transparency of what was actually perturbed. If corruption procedures are underspecified, it becomes difficult to interpret why a metric changed. For this reason, our pipeline emits realized diagnostics at row, column, and cell levels for each corruption setting. These diagnostics support quality control of the perturbation process itself, allowing readers to distinguish ``weak effect because corruption was harmless'' from ``weak effect despite substantial corruption realization''.

\subsection{Probability Calibration and Reliability Metrics}
Probability calibration has a long history in machine learning and remains central to decision systems. Classical and modern work has shown that good ranking performance does not guarantee calibrated probabilities \cite{platt2000probabilities,zadrozny2002transforming,niculescu2005predicting,naeini2015bbq,guo2017calibration,kull2017beta}. This distinction is operationally important when probabilities are interpreted as risk estimates or when fixed thresholds are used to trigger downstream actions.

Platt scaling introduced a parametric approach based on logistic mapping of classifier scores \cite{platt2000probabilities}. Subsequent work extended calibration strategies, including multiclass probability correction \cite{zadrozny2002transforming}, broad empirical evaluation across supervised learners \cite{niculescu2005predicting}, Bayesian binning \cite{naeini2015bbq}, and analyses of calibration behavior in modern deep models with strong results for temperature scaling \cite{guo2017calibration}. Beta calibration further generalized sigmoid-style mappings and has shown robust behavior in binary settings \cite{kull2017beta}.

The present study does not attempt to resolve calibration method superiority in a universal sense. Instead, it evaluates calibrators in a stress-testing context where input corruption shifts the probability distribution and potentially invalidates calibration fitted on clean validation data. This setup is important for deployment: calibration is often trained once and then reused under changing data quality conditions. A calibrator that performs well in clean settings may underperform or even worsen secondary metrics under severe corruption.

ECE is used here as a primary calibration-quality endpoint because it provides an interpretable summary of confidence-accuracy discrepancy. In addition, we retain log-loss and Brier-style perspectives in secondary reporting to avoid single-metric overinterpretation. This multi-metric approach is important because different calibration diagnostics can prioritize different failure patterns, especially under class imbalance and distributional stress.

\subsection{Threshold Governance and Decision Cost}
Threshold selection is often treated as a post-processing detail, but in many operational systems it is the direct control that governs intervention rates, alert load, and downstream resource allocation. A fixed threshold such as $0.5$ is convenient but can be misaligned with class prevalence, calibration drift, and asymmetric error costs. In practice, organizations frequently operate with explicit or implicit cost asymmetry, making threshold governance a policy question as much as a modeling question.

This manuscript formalizes threshold governance with three operating policies: fixed $0.5$, validation-tuned F1 threshold, and cost-sensitive threshold derived from error-cost ratio. The intent is to compare robustness not only of the predictive model but also of policy choices under corruption. This framing reflects real deployment decisions, where operators can often adjust thresholds faster than they can retrain models.

The policy analysis in this work is intentionally tied to expected cost. F1 and balanced accuracy provide useful performance summaries, but expected cost provides an interpretable link between metric behavior and operational consequence under explicit cost assumptions. This linkage allows reliability results to be discussed in terms that are actionable for stakeholders responsible for intervention budgets and service-level objectives.

\subsection{Positioning of This Study}
Relative to prior work, this study contributes an integrated benchmark perspective. The novelty is not a new calibrator, classifier, or significance test in isolation; rather, it is the coupling of structured operational corruptions, calibration-aware evaluation, policy-level analysis, and multiplicity-aware inference in one reproducible workflow.

Compared with dataset-shift and concept-drift literature \cite{quionero2009datasetshift,moreno2012datasetshift,gama2014conceptdrift}, our unit of analysis is explicitly \emph{pipeline fault families} (C1--C4) with realized corruption diagnostics and fixed severe anchors. Compared with calibration studies \cite{platt2000probabilities,guo2017calibration,kull2017beta}, calibration is not evaluated only on clean holdout data but under corruption-conditioned stress, and is linked to downstream threshold-policy utility. Compared with technical-debt/system-reliability framing \cite{sculley2015hidden}, this work adds a concrete, machine-auditable artifact chain (raw $\rightarrow$ summary $\rightarrow$ manuscript tables/figures) for peer verification.

This positioning motivates conservative claim language. Confirmatory language is restricted to the predefined primary family with explicit dataset-count bounds; broader setting-level patterns are treated as secondary or exploratory. The objective is an auditable applied reliability procedure rather than universal model-superiority claims.

\section{Methods}
\subsection{Research Goal and Questions}
The goal of this study is to evaluate the reliability of tabular machine-learning systems under structured operational data corruptions and to assess whether calibration and threshold-policy controls can reduce deployment risk.

\subsection{Research Method and Procedure}
We use a reproducible benchmark procedure with predefined datasets, models, corruption operators, severity grids, and inferential outputs. The design supports artifact-level traceability from raw experimental outputs to manuscript tables and figures.

\subsection{Benchmark Design}
The benchmark includes 12 OpenML binary datasets spanning socioeconomic, financial, industrial, medical, telecom, transportation, and justice-related domains \cite{vanschoren2014openml}. Table~\ref{tab:design} summarizes the global design.

\begin{table*}[!t]
\caption{Experimental design summary.\label{tab:design}}
\begin{tabularx}{\textwidth}{lX}
\toprule
\textbf{Component} & \textbf{Configuration} \\
\midrule
Datasets & 12 binary tabular datasets, 562,660 total rows \\
Model families & Logistic regression, random forest, histogram gradient boosting \\
Encoding policies & \texttt{ignore\_unknown}, \texttt{unknown\_bucket} \\
Seeds & 20 (0--19), stratified 70/15/15 train/validation/test split \\
Single corruptions & C1, C2, C3, C4 with severity $\alpha\in\{0.0,0.05,0.1,0.2,0.4\}$ \\
Compound corruption & C2+C4 with $(\alpha_{C2},\alpha_{C4})\in\{0.0,0.1,0.2,0.4\}\times\{0.0,0.2,0.4\}$ \\
Primary endpoints & AUC, F1@0.5, ECE \\
Additional endpoints & PR-AUC, balanced accuracy, Brier score \cite{brier1950verification}, log-loss, calibration diagnostics, policy cost \\
\bottomrule
\end{tabularx}
\end{table*}

The design intentionally combines breadth and controllability. Breadth is obtained through 12 datasets with heterogeneous prevalence, feature dimensionality, and missingness baselines. Controllability is obtained through fixed seeds, explicit severity grids, and deterministic corruption routines. This combination allows us to evaluate whether observed degradation patterns are robust across contexts while preserving enough structure for reproducible hypothesis testing.

All corruption experiments are applied to test partitions only; train and validation splits remain clean within each seed. This setup reflects a common deployment asymmetry in which model training can be quality-controlled, while inference-time data quality is less predictable. By preserving clean train/validation data, the benchmark isolates the effect of test-time data degradation and avoids conflating corruption effects with training instability.

Severity settings are defined as an intervention intensity rather than a direct guarantee of exact affected-cell percentage. Realized corruption impact is therefore measured and logged through diagnostics, including row impact rate, column impact rate, and cell-level intervention rate. This is important because dataset schemas differ substantially; a nominally identical severity parameter can correspond to different absolute perturbation volumes depending on feature count and type composition.

The severe-versus-baseline contrast in this manuscript is anchored at 0.4 versus 0.0 for single corruptions and $(0.4,0.4)$ versus $(0.0,0.0)$ for the compound case. This severe anchor was chosen for interpretation consistency across families and aligns with the benchmark objective of stress testing rather than mild sensitivity analysis.

\subsection{Dataset Profiles and Experimental Setup}
Table~\ref{tab:dataset_profiles} reports the dataset profile summary used for this manuscript, including OpenML IDs, sample size, feature composition, class prevalence, and maximum column missingness.

\begin{table*}[!t]
\caption{Dataset profile summary (from \texttt{../results/tables/table\_dataset\_profiles.csv}).\label{tab:dataset_profiles}}
\footnotesize
\fittableinput{../results/tables/table_dataset_profiles.tex}
\end{table*}

Dataset inclusion followed predefined protocol DSP-20260225-01 (see repository documentation). Eligibility checks covered binary-target validity, minimum row and feature thresholds, prevalence bounds, leakage-proxy screening, and pinned OpenML provenance. Three datasets (\texttt{airlines}, \texttt{diabetes130us}, and \texttt{aps\_failure}) were retained under explicit protocol waivers to preserve transportation, clinical, and industrial coverage in the benchmark.

Models were implemented with scikit-learn \cite{pedregosa2011scikit}, with random forest \cite{breiman2001random}, gradient boosting foundations \cite{friedman2001gbm}, and histogram gradient boosting via the official scikit-learn implementation \cite{sklearnHGBDocs}. The benchmark uses fixed random seeds, two categorical handling policies, and consistent train/validation/test partitioning across corruption settings.

Dataset sourcing follows pinned OpenML identifiers to improve reproducibility and reduce ambiguity from dataset-name and version drift. For each included dataset, the submission artifacts retain source identifiers, row counts, feature composition, prevalence, missingness diagnostics, and deterministic local hash manifests for the archived dataset snapshots.
OpenML pinning in this study is at the dataset record and version level rather than an external content-hash level for remote storage. Consequently, OpenML IDs and versions act as deterministic protocol anchors, while byte-level immutability is guaranteed only for the archived local submission artifacts.
Dataset-level provenance is anchored to OpenML data records for all 12 datasets:
\cite{openml_d1590_adult,openml_d42493_airlines,openml_d41138_apsfailure,openml_d1461_bank_marketing,openml_d44162_compass,openml_d42477_credit_default,openml_d43903_diabetes_hospitals_fairlearn,openml_d45022_diabetes130us,openml_d151_electricity,openml_d41162_kick,openml_d43904_law_school_admission,openml_d45568_telco_churn}.

The split protocol is stratified by target class for each seed. Each run uses a 70/15/15 split across train, validation, and test partitions. Stratification helps stabilize threshold and calibration fitting under class imbalance. Twenty seeds (0--19) are used to estimate within-setting variability while preserving feasible full-pipeline compute.

Preprocessing uses a column-wise pipeline with median imputation and standardization for numeric features and most-frequent imputation for categorical features. Two encoding policies are evaluated to reflect common operational choices for unseen categories. Under \texttt{ignore\_unknown}, unseen categories at inference are ignored by one-hot encoding. Under \texttt{unknown\_bucket}, unseen categories are mapped to an explicit \texttt{UNK} bucket before one-hot transformation. This comparison is particularly relevant under C3 categorical drift, where previously unseen tokens are injected by design.

Model families were selected for deployment relevance in tabular workflows. Logistic regression provides a linear baseline with high interpretability and stable optimization. Random forest adds non-linear partitioning and interaction capture through bagged trees \cite{breiman2001random}. Histogram gradient boosting provides boosted tree capacity with histogram-based efficiency and native support for probabilistic binary classification under log-loss optimization \cite{friedman2001gbm,sklearnHGBDocs}.

Default model settings were not aggressively hyper-optimized for each dataset. This was intentional. The objective of this benchmark is to compare reliability behavior across corruption conditions using stable model templates, not to maximize clean benchmark leaderboards. Excessive per-dataset tuning could obscure corruption-induced effects by introducing additional model-selection variance and reduce portability of findings to practical monitoring contexts.

Threshold-policy analysis reuses validation-derived thresholds at test time, mirroring common operational deployment where threshold updates are periodic and not continuously re-estimated on live labels. This design makes policy-level conclusions more realistic for decision systems with delayed outcome feedback.

\begin{table*}[!t]
\caption{Implementation parameterization used for all manuscript artifacts.\label{tab:implparams}}
\begin{tabularx}{\textwidth}{lX}
\toprule
\textbf{Component} & \textbf{Exact implementation setting} \\
\midrule
Data split & Stratified \texttt{train\_test\_split}: 70/15/15 with seeds \{0,\ldots,19\} \\
Preprocessing (numeric) & \texttt{SimpleImputer(strategy='median')} + \texttt{StandardScaler()} \\
Preprocessing (categorical) & \texttt{SimpleImputer(strategy='most\_frequent')} + one-hot encoding (\texttt{ignore\_unknown} or explicit \texttt{UNK} bucket) \\
Logistic regression & \texttt{solver='lbfgs'}, \texttt{max\_iter=1000}, default L2, \texttt{random\_state=seed} \\
Random forest & \texttt{n\_estimators=300}, \texttt{n\_jobs=-1}, \texttt{random\_state=seed} \\
Histogram gradient boosting & \texttt{max\_iter=200}, \texttt{learning\_rate=0.1}, \texttt{max\_depth=None}, \texttt{loss='log\_loss'}, \texttt{random\_state=seed} \\
C1 duplication volume & $k=\mathrm{round}(\alpha n)$ duplicated rows; sampling with replacement when anchor cohort size $<k$ \\
C1 anchor mode & Manuscript run uses \texttt{label\_agnostic}; optional \texttt{label\_informed} mode is implemented for sensitivity reruns (\texttt{C1\_ANCHOR\_MODE} env var) \\
C2 masking volume & row count $k=\mathrm{round}(\alpha n)$; masked-column fraction $\max(\alpha,1/|\mathcal{C}|)$ \\
C3 token generation & token families \{\texttt{new}, \texttt{legacy}, \texttt{vendor}, \texttt{hash}\}; family count $\min(4,\max(2,\mathrm{round}(2+2\alpha)))$; token count $\max(2,\mathrm{round}(2+6\alpha))$ \\
C4 numeric variants & per-column variant probabilities: scale-shift 0.35, shift-only 0.25, rounding 0.20, clipping 0.20 \\
C4 targeted-column fraction & $\min\!\left(1,\max\!\left(0.25 + 0.75\alpha,\;1/|\mathcal{N}|\right)\right)$ where $|\mathcal{N}|$ is numeric-column count \\
C4 guardrail & per-cell absolute delta cap $(0.25 + 2.00\alpha)\sigma_{\text{column}}$ \\
Corruption RNG seed offsets & C1: \texttt{seed+11}; C2: \texttt{seed+21}; C3: \texttt{seed+31}; C4: \texttt{seed+41}; C2+C4: base \texttt{seed+51} with C4 inner stream \texttt{+10000} \\
Inoperative-family fallback & If required column type is absent (for example C3 with no categorical columns), operator is a no-op and diagnostics record zero affected rows/cells \\
ECE & 15 equal-width bins over clipped probabilities $p \in [10^{-6}, 1 - 10^{-6}]$; first bin closed, others right-closed \\
\bottomrule
\end{tabularx}
\end{table*}

\subsection{Corruption Operators}
Let $(X,y)$ denote clean features and labels, and let $T_c(X,\alpha,\omega)$ denote corruption family $c$ with severity $\alpha$ and seeded stochastic stream $\omega$. Corrupted inputs are generated as:
\begin{equation}
X^{(c,\alpha)} = T_c(X,\alpha,\omega).
\label{eq:corr_operator}
\end{equation}

For compound corruption C2+C4, operators are applied sequentially:
\begin{equation}
T_{C2+C4}(X,\alpha_{2},\alpha_{4},\omega)
=
T_{C4}\!\left(T_{C2}(X,\alpha_{2},\omega_{2}),\alpha_{4},\omega_{4}\right).
\label{eq:compound_operator}
\end{equation}
Equations~\eqref{eq:corr_operator} and \eqref{eq:compound_operator} define the single and compound corruption transformations used in all experiments.

The corruption families are:
\begin{itemize}
\item \textbf{C1 (cohort-biased duplication):} duplicates $k=\mathrm{round}(\alpha n)$ rows from anchor cohorts to emulate join amplification.
\item \textbf{C2 (structured missingness):} masks anchor-selected values with severity-coupled row and column sampling.
\item \textbf{C3 (multi-pattern categorical drift):} injects unseen tokens from deterministic \texttt{new}/\texttt{legacy}/\texttt{vendor}/\texttt{hash} token families.
\item \textbf{C4 (multi-variant numerical inconsistency):} applies scale-shift, shift-only, rounding, and clipping perturbations to selected numeric columns.
\item \textbf{C2+C4:} sequential composition of C2 then C4 with independent streams.
\end{itemize}

\subsubsection{C1: Cohort-Biased Duplication}
C1 emulates join-induced row amplification. A severity-coupled number of rows is duplicated from anchor cohorts selected through frequency-based categorical groups when available, with numeric fallback when needed. The default manuscript configuration uses \emph{label-agnostic} anchoring. A label-informed stress-synthesis mode is implemented (\texttt{C1\_ANCHOR\_MODE='label\_informed'}) and retained for sensitivity analysis only.

C1 primarily perturbs effective sample weighting and local class composition in affected cohorts. In many settings, this can alter threshold-sensitive metrics more strongly than ranking metrics, because duplicated observations reinforce certain score regions without necessarily reordering the entire score distribution. Label-informed C1 findings are interpreted as descriptive stress-synthesis sensitivity evidence and are excluded from confirmatory ranking claims.

\subsubsection{C2: Structured Missingness}
C2 emulates key-mismatch and field-dropout behavior. It selects severity-coupled row subsets, preferentially from anchored cohorts, and masks a severity-coupled fraction of columns. The masking strategy is structured rather than independent cell dropout, creating correlated missingness patterns that are common in pipeline faults.

At low severity, C2 may mask a small number of columns over a limited row subset. At high severity, both row coverage and column coverage increase, creating large coherence blocks of missing information. This design is operationally important because model sensitivity to coordinated missingness can differ substantially from sensitivity to random single-cell omissions.

\subsubsection{C3: Multi-Pattern Categorical Drift}
C3 injects unseen categorical tokens from multiple families (\texttt{new}, \texttt{legacy}, \texttt{vendor}, \texttt{hash}) with severity-coupled token diversity. The objective is to emulate realistic schema and value drift where category novelty is not monolithic. In production systems, new categories may come from vendor updates, code-system revisions, hashed surrogates, or legacy fallback values.

By varying token families and column subsets, C3 stresses both encoding policy behavior and model robustness to out-of-vocabulary conditions. This is especially relevant when deployments rely on one-hot encoding with limited unknown-handling semantics. Hash-family token identifiers are generated with a process-stable SHA-256 mapping from $(\text{column},\text{seed},\text{token-index})$, avoiding Python runtime hash randomization and preserving determinism under fixed seeds.

\subsubsection{C4: Multi-Variant Numerical Inconsistency}
C4 applies one perturbation variant per selected numeric column from a set including scale-shift, shift-only, rounding, and clipping interventions. Variants are chosen to mimic common measurement inconsistency pathways, such as hidden unit changes, numeric truncation, and integration clipping rules.

Unlike pure additive noise, C4 can produce distributional distortions with structural asymmetry. For example, clipping can disproportionately affect tails, while rounding can reduce effective resolution. These distortions may impact calibration and threshold behavior even when AUC changes are moderate.

\subsubsection{Compound C2+C4}
Compound corruption applies C2 first and C4 second with independent random streams. The ordering is intentional: missingness is introduced before numerical perturbation, simulating data-loss plus measurement inconsistency in the same downstream batch. This compound setting provides a stress proxy for multi-fault scenarios that are frequently more harmful than isolated single-fault events.

Compound diagnostics are emitted with separate C2 and C4 realized-parameter traces, and overall row/column impact is computed on the union of changed cells in the composed transformation, enabling decomposition without overstating aggregate realization.

For metric $\phi$, severe-minus-baseline effect at setting $s$ is:
\begin{equation}
\Delta\phi_s^{(c)}=\phi_s(\alpha_h)-\phi_s(\alpha_0),
\label{eq:delta_metric}
\end{equation}
where $\alpha_h=0.4$ and $\alpha_0=0.0$ for single families; for compound corruption, $(0.4,0.4)$ is compared against $(0.0,0.0)$.
All severe-minus-baseline results in Section~4 use Eq.~\eqref{eq:delta_metric}.

\subsection{Metrics and Endpoint Interpretation}
Primary endpoints are AUC, \textbf{F1@0.5}, and ECE. This triad was selected to span ranking quality, fixed-operating-point classification quality, and probability reliability. The benchmark also stores PR-AUC, balanced accuracy, Brier score, log-loss, adaptive ECE, calibration slope/intercept, and reliability-gap diagnostics for supportive interpretation.

AUC captures rank discrimination across all thresholds and is invariant to monotonic score transformations. F1@0.5 captures behavior at a fixed deployment baseline threshold and is therefore sensitive to prevalence context. Threshold-optimized F1 is reported separately in the policy section. ECE captures the aggregate discrepancy between predicted confidence and empirical event frequency under binning. Joint interpretation is necessary because each metric captures a different reliability axis.

Secondary diagnostics support failure-mode interpretation. For example, calibration slope and intercept help identify over- and under-confidence tendencies beyond scalar ECE. Reliability-gap summaries (mean and maximum) indicate whether calibration errors are diffuse across bins or concentrated in high-risk score regions. Expected decision cost translates classification outcomes to operational impact under explicit cost assumptions. In implementation, ECE uses $B=15$ equal-width bins on clipped probabilities $p \in [10^{-6}, 1 - 10^{-6}]$.

The severe-versus-baseline delta formulation standardizes interpretation across heterogeneous datasets. Absolute metric values can vary strongly with base task difficulty; delta effects provide a more direct estimate of reliability loss under corruption stress.

\subsection{Calibration and Threshold Governance}
Evaluated calibrators were uncalibrated baseline, temperature scaling, Platt scaling, isotonic regression, and beta calibration \cite{platt2000probabilities,zadrozny2002transforming,niculescu2005predicting,guo2017calibration,kull2017beta}.

Temperature scaling optimizes validation negative log-loss:
\begin{equation}
\begin{aligned}
\hat{p}^{(T)}_i
&=
\sigma\!\left(\frac{\operatorname{logit}(\hat{p}_i)}{T}\right), \\
T
&=
\arg\min_{T>0}
\left[
-\frac{1}{n_v}\sum_{i=1}^{n_v}
\left(
y_i\log \hat{p}^{(T)}_i
\right.\right. \\
&\qquad\qquad\left.\left.
+ (1-y_i)\log\!\left(1-\hat{p}^{(T)}_i\right)
\right)
\right].
\end{aligned}
\label{eq:temperature}
\end{equation}

Calibration quality is reported by expected calibration error (ECE) \cite{naeini2015bbq,guo2017calibration}:
\begin{equation}
\mathrm{ECE}
=
\sum_{b=1}^{B}\frac{|I_b|}{N}
\left|
\mathrm{acc}(I_b)-\mathrm{conf}(I_b)
\right|,
\label{eq:ece}
\end{equation}
where $I_b$ is bin $b$, $\mathrm{acc}(I_b)$ is the empirical observed outcome rate in the bin, and $\mathrm{conf}(I_b)$ is mean predicted probability.
For binary outcomes, $\mathrm{acc}(I_b)$ is implemented as the empirical positive-class rate $\frac{1}{|I_b|}\sum_{i\in I_b} y_i$ (not hard-label classification accuracy), and $\mathrm{conf}(I_b)$ is the mean predicted probability.
In code, $B=15$ equal-width bins are used; the first bin is closed on both ends and subsequent bins are right-closed, matching the implementation used for all reported artifacts.

Policy evaluation used three thresholds per variant: fixed $t=0.5$, validation-tuned F1 threshold, and a cost-sensitive threshold
\begin{equation}
t^*=\frac{C_{FP}}{C_{FP}+C_{FN}},
\label{eq:cost_threshold}
\end{equation}
with $C_{FP}=1$, $C_{FN}=5$, yielding $t^*=1/6$.

The validation operating point was selected by:
\begin{equation}
\hat{t}_{\mathrm{F1}}
=
\arg\max_{t\in\mathcal{G}} \mathrm{F1}(t),
\qquad
\mathcal{G}=\{0.01,0.02,\ldots,0.99\}.
\label{eq:f1_threshold}
\end{equation}

Expected decision cost is:
\begin{equation}
\mathcal{C}(t)=\frac{C_{FP}\cdot FP(t)+C_{FN}\cdot FN(t)}{N}.
\label{eq:expected_cost}
\end{equation}
Threshold governance and expected cost are operationalized with Eqs.~\eqref{eq:cost_threshold}--\eqref{eq:expected_cost}.

Temperature scaling in this benchmark is fitted by minimizing validation negative log-loss over a single positive scale parameter. Platt scaling fits a logistic model over logit-transformed base scores. Isotonic regression provides a non-parametric monotonic mapping, and beta calibration uses log-transformed probability features $(\log p,\log(1-p))$ with the same logistic optimizer and regularization configuration \cite{kull2017beta}. Each calibrator is trained on clean validation predictions within a seed-model-encoding setting and then applied to corresponding test predictions across corruption conditions.

This calibration protocol reflects a realistic deployment practice: calibrators are often estimated on reference validation data and reused until explicit recalibration cycles occur. The benchmark therefore evaluates not only ``best possible calibration under each corruption'' but also ``deployed calibration persistence under stress'', which is more relevant for operational reliability audits.

Threshold governance compares three policies:
\begin{itemize}
\item \textbf{Fixed 0.5:} conventional default threshold.
\item \textbf{Validation-tuned F1:} threshold maximizing validation F1 over a discrete grid.
\item \textbf{Cost-sensitive:} threshold implied by error-cost asymmetry.
\end{itemize}
The policy comparison is performed for both uncalibrated and temperature-scaled variants in the underlying artifacts, while the manuscript emphasizes uncalibrated comparisons for direct interpretation of corruption-driven policy benefit.

By reporting both F1 deltas and expected-cost deltas against a fixed-threshold baseline, the analysis distinguishes between ``threshold improves balance-oriented metric'' and ``threshold improves operational objective under explicit costs''. This distinction is important in applications where intervention capacity and miss-cost asymmetry are primary constraints.

\subsection{Statistical Inference}
Inference is reported at two scopes. First, \emph{setting-level} severe-vs-baseline paired seed deltas were analyzed with paired $t$-tests, Wilcoxon signed-rank tests \cite{wilcoxon1945individual}, bootstrap 95\% confidence intervals, and paired effect size $d_z$. Wilcoxon was computed on non-zero paired deltas (\texttt{zero\_method='wilcox'}), and both total paired count and non-zero effective paired count are logged per setting.

Second, and treated as the primary inferential unit for manuscript conclusions, a \emph{dataset-level} paired estimand was computed for the predefined family (uncalibrated C2 and C2+C4; AUC/F1/ECE). For each dataset and seed, severe-minus-baseline deltas were averaged over model family and encoding policy, then tested across paired seeds ($n=20$ per dataset). This reduces dependence from reusing many correlated setting-level tests as stand-alone evidence.

Multiplicity correction used Benjamini--Hochberg. Setting-level analyses were corrected within each hypothesis family and metric, while dataset-level predefined analyses were corrected within metric across datasets:
\begin{equation}
q_{(i)}=\min_{j\ge i}\left(\frac{m}{j}p_{(j)}\right).
\label{eq:bh}
\end{equation}
For pairwise policy and calibration batteries, BH correction is applied once over the full battery in each table (all listed scenarios and metrics), and reported q-values are interpreted within that battery-wide family.
Pairwise calibration and policy comparisons involve multiple contrasts over n=12 dataset clusters. These rankings should therefore be interpreted as exploratory rather than confirmatory.

The predefined primary family remains uncalibrated C2 and C2+C4 for AUC/F1/ECE. With 20 paired seeds per dataset, setting-level uncertainty is reduced, while cross-dataset confirmatory scope remains bounded by the number of datasets ($n=12$). Accordingly, p-values and BH counts are interpreted together with effect magnitudes and confidence intervals at dataset and cross-dataset summary levels. Secondary and exploratory analyses are explicitly separated to limit post hoc interpretation bias.
Using a two-sided paired-$t$ approximation at $\alpha=0.05$ with target power $0.80$ and $n=12$ dataset-level units, the minimum detectable standardized paired effect is approximately $d_z\approx 0.89$; smaller true effects can therefore be directionally informative but underpowered for confirmatory significance.

Statistical assumptions are treated conservatively. Bootstrap confidence intervals in this work resample the analysis unit used by each table (settings for setting-level summaries, datasets for cross-dataset summaries), but they do not remove all dependence induced by the nested dataset/model/encoding design. Therefore, bootstrap CIs are interpreted as uncertainty summaries under the stated design, not as guarantees under full hierarchical independence.

To check sensitivity to nested dependence, we additionally computed a cluster-aware bootstrap that resamples datasets, then model and encoding settings within dataset, and finally seeds within setting. This hierarchical sensitivity is treated as a robustness check for headline direction and CI support, not as a replacement for the predefined primary analysis. For policy and calibration comparisons, inference is summarized on dataset-cluster means to avoid treating correlated within-dataset settings as independent observations. Severe-summary uncertainty intervals for policy and calibration are likewise reported with dataset-cluster bootstrap CIs ($n=12$ clusters).

For the predefined six primary hypotheses, direct $p$/$q$ reporting is provided in Table~\ref{tab:primaryhyp} to complement count-based summaries.

\subsection{Notation}
For clarity: $c$ denotes corruption family; $\alpha$ denotes corruption severity; $\omega$ denotes seeded stochastic stream; $\phi$ denotes a performance metric; $\Delta\phi$ denotes severe-minus-baseline effect; $t$ denotes decision threshold; and $\mathcal{C}(t)$ denotes expected decision cost. Equations \eqref{eq:corr_operator}--\eqref{eq:bh} define the formal analysis pipeline.

\section{Results}
\subsection{Cross-Dataset Effects}
Table~\ref{tab:cross} reports cross-dataset mean performance deltas and bootstrap confidence intervals under severe corruptions. In this run, C2 has the largest mean harmful single-family effect, while C2+C4 has the largest mean harmful compound effect on AUC and F1. Figure~\ref{fig:crosssummary} provides the visual summary.

The estimand in Table~\ref{tab:cross} is an unweighted mean across datasets of dataset-level severe-minus-baseline deltas, where each dataset-level delta is itself an unweighted mean across encoding-policy and model-family settings. Among single corruptions, C2 produces the largest average degradation in both discrimination and thresholded classification quality. The mean $\Delta$AUC under C2 is -0.031 and the mean $\Delta$F1 is -0.055, both with confidence intervals that remain on the harmful side.

For C3, eligibility is restricted to datasets with at least one categorical predictor (\texttt{n\_cat\_cols > 0}), because C3 is otherwise inoperative by design. Under this eligibility-aware aggregation (10 datasets), C3 shows larger mean harmful shifts than C4 on AUC and F1 in this run ($\Delta$AUC -0.013, $\Delta$F1 -0.023), while C4 remains weaker but still generally harmful.

C1 behaves differently. Its manuscript-default (label-agnostic) mean AUC effect is slightly negative, while mean F1 degrades only slightly and mean ECE increases. This is consistent with the design of C1 as a cohort-weighting perturbation. Duplication can alter class balance and threshold behavior in affected regions without necessarily causing strong global rank deterioration.

The combined C2+C4 condition yields larger degradation than either corruption family alone for AUC and F1. The mean $\Delta$AUC is -0.040 and mean $\Delta$F1 is -0.069.

ECE behavior is more nuanced. C2 increases ECE by approximately +0.012 on average and C2+C4 by +0.016, while C3 and C4 show smaller average increases. The confidence interval for C4 ECE includes values close to zero, indicating that numerical inconsistency does not uniformly increase calibration error across all datasets at the same magnitude as C2.

Row-weighted sensitivity analysis preserves the same qualitative ordering for AUC and F1, indicating that conclusions are not an artifact of equal-dataset weighting alone.

Overall, the cross-dataset aggregate indicates a stable descriptive ordering in this benchmark among non-label-informed single-family channels: C2 has the largest mean harmful single-family shifts, C3 is typically next among categorical-eligible datasets, and C4 is generally weaker but non-negligible; compound C2+C4 remains most harmful on mean AUC/F1 deltas.

\begin{table*}[!t]
\caption{Cross-dataset severe effects (0.4 minus 0.0), from \texttt{../results/tables/table\_cross\_dataset\_effects.csv}. Estimand: dataset-unweighted mean of dataset-level severe-minus-baseline deltas, where each dataset-level delta averages over model and encoding settings. Uncertainty: bootstrap 95\% CI over datasets. Eligibility: C1/C2/C4/C2+C4 use all 12 datasets; C3 uses categorical-eligible datasets only (\texttt{n\_cat\_cols > 0}, $n=10$). Harmful direction is $\Delta$AUC/$\Delta$F1 $< 0$ and $\Delta$ECE $> 0$.\label{tab:cross}}
\small
\fittableinput{../results/tables/table_cross_dataset_effects.tex}
\end{table*}

\subsection{Primary Effect Diagnostics}
Primary-family diagnostics are summarized in Table~\ref{tab:primaryhyp}, which reports direct $p$/$q$ values for the six predefined primary hypotheses (C2/C2+C4 $\times$ AUC/F1/ECE) from cross-dataset tests over dataset-level deltas.

Research Question 2 asks whether severe compound corruption is stronger than severe single-family corruption. At both setting and dataset levels, C2+C4 is more harmful than C2 for AUC and F1 in this run. This interpretation is anchored to effect-size magnitude, harmful-direction consistency, and bootstrap CIs; test statistics are treated as secondary diagnostics.

Given $n=20$ paired seeds per dataset, inferential stability is improved at the seed level, while cross-dataset inference remains bounded by $n=12$ datasets. Primary interpretation therefore emphasizes agreement between effect-size magnitude, harmful-direction consistency across datasets, and the direct $p$/$q$ values in Table~\ref{tab:primaryhyp}. All primary hypotheses are supported in the harmful direction for the primary metrics after BH correction.

\begin{table*}[!t]
\caption{Direct predefined-hypothesis inference table (\texttt{../results/tables/table\_primary\_hypothesis\_tests.csv}). Each row is one predefined primary hypothesis (C2/C2+C4 $\times$ AUC/F1/ECE). Deltas are dataset-level severe-minus-baseline effects; two-sided one-sample $t$ and Wilcoxon tests are reported with BH correction across the six tests. Harmful direction is $\Delta$AUC/$\Delta$F1 $< 0$ and $\Delta$ECE $> 0$; inference is dataset-cluster-aware by construction.\label{tab:primaryhyp}}
\scriptsize
\fittableinput{../results/tables/table_primary_hypothesis_tests.tex}
\end{table*}


\subsection{Heterogeneity}
C2 effects are heterogeneous across datasets. The largest C2 F1 degradations are observed in \texttt{credit\_default} (-0.242), \texttt{aps\_failure} (-0.116), and \texttt{bank\_marketing} (-0.062), while \texttt{airlines} shows an opposite-direction case (+0.022). This pattern supports dataset- and pipeline-specific sensitivity analysis instead of single-number robustness claims.

Under severe C2 (uncalibrated), model-family means are: \texttt{hist\_gb} ($\Delta$AUC -0.031, $\Delta$F1 -0.058, $\Delta$ECE +0.013), \texttt{logreg} ($\Delta$AUC -0.034, $\Delta$F1 -0.050, $\Delta$ECE +0.018), and \texttt{random\_forest} ($\Delta$AUC -0.028, $\Delta$F1 -0.056, $\Delta$ECE +0.005). No model dominates all endpoints.

Heterogeneity is visible not only in magnitude but also in direction for selected settings. For example, one dataset can show moderate F1 recovery under a corruption that is harmful on average across datasets. This does not invalidate cross-dataset means; instead, it indicates that context-specific data geometry and class structure influence how each corruption family manifests at metric level.

A setting-level sign analysis reinforces this point. Under severe C2, AUC degradation is harmful in all datasets, but F1 and ECE include isolated opposite-direction cases. Under C3 and C4, harmful direction counts are lower than C2, consistent with a weaker and less uniform stress profile. C1 has the highest directional variability, matching its cohort-weighting mechanism.

Model-family comparisons under severe C2+C4 show the same qualitative pattern as C2 but with larger average losses: \texttt{hist\_gb} ($\Delta$AUC -0.044, $\Delta$F1 -0.075, $\Delta$ECE +0.020), \texttt{logreg} ($\Delta$AUC -0.040, $\Delta$F1 -0.062, $\Delta$ECE +0.022), and \texttt{random\_forest} ($\Delta$AUC -0.037, $\Delta$F1 -0.069, $\Delta$ECE +0.007). Again, no family is uniformly best over all reliability axes.

Encoding-policy heterogeneity is comparatively small at benchmark level for severe C3 in this run. The mean unknown-minus-ignore effect is close to zero for AUC and ECE (about 0.000 for both) and slightly negative for F1 (about -0.001), but with non-trivial setting-level spread. This indicates that choosing an unknown-handling policy can still matter locally, especially in datasets where category novelty affects high-impact cohorts.
Two datasets in this benchmark (\texttt{diabetes130us}, \texttt{aps\_failure}) have no categorical predictors, so C3 is inoperative there; C3 family-level aggregates are therefore reported on categorical-eligible datasets only.

\subsection{Calibration Under Corruption}
Against uncalibrated clean baseline (ECE 0.030, log-loss 0.407), all calibrators reduce ECE. For severe single-corruption scenarios (C2/C4 over all datasets; C3 over categorical-eligible datasets), $\Delta$ECE versus uncalibrated is: temperature scaling -0.008 (95\% CI [-0.017, -0.002]), Platt -0.010 ([-0.018, -0.004]), isotonic -0.010 ([-0.019, -0.004]), and beta -0.011 ([-0.019, -0.004]).

Research Question 4 focuses on calibrator stability under stress rather than clean-only gains. On clean data, all calibrators improve ECE relative to uncalibrated baseline. Under severe corruption, the pattern remains beneficial for ECE across calibrators on average, but log-loss behavior differentiates methods.

Beta calibration shows the strongest average improvement in ECE, followed closely by isotonic regression and Platt scaling, while temperature scaling provides smaller but still consistent improvements. For log-loss, beta and Platt show the largest average improvements in this run. However, single-severe ECE CIs for temperature/Platt/isotonic/beta overlap substantially, so fine-grained ranking remains uncertain without formal pairwise testing. Isotonic remains competitive on ECE but its single-severe log-loss delta has a confidence interval spanning zero (-0.002, 95\% CI [-0.010, +0.005]), indicating weaker directional certainty for log-loss improvement.

Under severe compound C2+C4, calibrator gains contract relative to single-family severe scenarios, which is expected under stronger distributional stress. ECE deltas remain negative on average for all calibrators; compound log-loss CIs span zero for each calibrator in this run.

A practical interpretation is that calibration should be evaluated as a stress-conditioned control, not a one-time clean optimization step. The benchmark results show that ``calibration works'' is too coarse; calibrator choice and stress context both matter for downstream reliability, and fine-grained calibrator ordering should be treated as descriptive rather than definitive.

\begin{table*}[!t]
\caption{Calibration deltas versus uncalibrated baseline with dataset-cluster bootstrap 95\% CIs (\(n=12\) clusters), from \texttt{../results/tables/table\_calibration\_comparison\_severe.csv}. Single-severe rows pool C2/C4 on all datasets and C3 on categorical-eligible datasets only; compound rows use C2+C4 at (0.4,0.4). The table reports explicit cluster counts (\(n_{\mathrm{single,cl}}\), \(n_{\mathrm{compound,cl}}\)) per calibrator. Negative $\Delta$ECE / $\Delta$ log-loss indicates improvement versus uncalibrated baseline.\label{tab:calibtable}}
\small
\fittableinput{../results/tables/table_calibration_comparison_severe.tex}
\end{table*}

\subsection{Policy-Level Utility}
For uncalibrated severe single scenarios (C2/C4 over all datasets; C3 over categorical-eligible datasets), validation-tuned thresholds improve F1 by +0.090 (95\% CI [+0.047, +0.140]) and reduce expected cost by -0.223 ([-0.376, -0.100]) versus fixed 0.5. Cost-sensitive thresholds reduce expected cost by -0.288 ([-0.438, -0.158]). Under severe compound C2+C4, cost-sensitive policy reduces expected cost by -0.322 ([-0.476, -0.184]) relative to fixed threshold.

Research Question 5 asks whether threshold governance can recover utility under corruption. The answer is affirmative in this benchmark. Both adaptive policies substantially outperform fixed 0.5 in severe conditions when evaluated against expected cost, and validation-tuned thresholds additionally provide strong F1 recovery.

Validation tuning produces the largest F1 improvement, while the cost-sensitive threshold yields the largest reduction in expected decision cost. This trade-off is stable across the severe scenarios evaluated in the benchmark.

The same qualitative pattern persists in severe compound scenarios, where policy adaptation is even more consequential because baseline fixed-threshold performance is lower. Any strict ordering beyond this descriptive pattern should therefore be read together with multiplicity-adjusted pairwise test outputs.

From an operational governance viewpoint, these results support explicit threshold review cycles tied to data-quality monitoring. A static threshold after corruption-induced score shift can impose avoidable cost. Policy adaptation is not a replacement for data-quality remediation, but it is a practical risk-reduction lever when remediation latency is non-zero.

\begin{table*}[!t]
\caption{Threshold-policy comparison for uncalibrated predictions with dataset-cluster bootstrap 95\% CIs (\(n=12\) clusters), from \texttt{../results/tables/table\_threshold\_policy\_comparison\_severe.csv}. Single-severe rows pool C2/C4 on all datasets and C3 on categorical-eligible datasets only; compound rows use C2+C4 at (0.4,0.4). The table reports explicit cluster counts (\(n_{\mathrm{single,cl}}\), \(n_{\mathrm{compound,cl}}\)). Positive $\Delta$F1 and negative $\Delta$Cost are improvements versus fixed-threshold \(t=0.5\).\label{tab:policytable}}
\centering
\scriptsize
\fittableinput{../results/tables/table_threshold_policy_comparison_severe.tex}
\end{table*}

The same qualitative ordering persists under alternative error-cost ratios derived from the stored confusion matrices: adaptive policies continue to reduce expected cost relative to fixed-threshold deployment, and the cost-sensitive threshold remains the lowest-cost option for explicit cost minimization.

\subsection{Corruption Realization Diagnostics}
The analysis above reports model-level outcomes; this subsection summarizes whether corruption operators realized the intended stress intensity. Aggregated diagnostics show that severe single-corruption settings affected roughly 38.3\% of rows on average across families. Mean severe single column-impact rate was approximately 22.4\%, and mean severe single cell-impact rate was approximately 9.0\%.

By family, realized severe row impact is near 40\% for C1, C2, and C4, and slightly lower for C3 due to family-specific column eligibility behavior. For C2 severe settings, mean severe column impact is approximately 40.4\% and mean severe cell impact is approximately 16.2\%, consistent with broad structured missingness. For C4 severe settings, mean severe cell impact is approximately 12.0\%, with additional variability induced by numeric-variant selection.

Compound C2+C4 severe settings show stronger realized stress than single families, with mean row impact near 40\%, mean column impact near 43.9\%, and mean cell impact near 28.2\%. Realized C2 targeted-missingness rates remain close to complete masking in targeted cells (about 98.8\% newly missing among targeted cells), while C4 perturbation magnitude is controlled through standard-deviation-based guardrails in implementation diagnostics.

These diagnostics provide a consistency check for interpretation. Strong observed degradation under C2 and C2+C4 is not an artifact of weak perturbation realization; the diagnostics confirm substantial intervention intensity at severe settings.

\subsection{Integrated Answer to Research Questions}
Combining all result views yields a consistent narrative:
\begin{itemize}
\item \textbf{RQ1 (dominant families):} C2 shows the largest mean harmful single-family effects in this benchmark run; C3 is typically second on categorical-eligible datasets.
\item \textbf{RQ2 (compound strength):} C2+C4 exceeds single-family severe effects on average for AUC and F1.
\item \textbf{RQ3 (heterogeneity):} effect size and even effect direction vary across datasets/settings, requiring context-aware interpretation.
\item \textbf{RQ4 (calibration robustness):} calibration helps on average under stress; beta shows the largest mean severe ECE improvement, with isotonic and Platt close behind and overlapping CIs across several calibrators.
\item \textbf{RQ5 (policy governance):} threshold adaptation delivers substantial utility recovery and cost reduction versus fixed-threshold deployment.
\end{itemize}

The integrated result is therefore not ``one model wins'' but ``governance stack matters''. Data-quality stress, calibration behavior, and threshold policy interact to determine deployment reliability.

\subsection{Figures}
The repository contains a larger figure set, while this manuscript retains only the global summary visuals most relevant to the journal article.

\begin{figure*}[!t]
\centering
\includegraphics[width=0.92\textwidth]{../results/plots/paper__cross_dataset__severe_deltas.png}
\caption{Cross-dataset severe-minus-baseline effects with bootstrap 95\% confidence intervals ($\Delta$ metric = severe minus baseline; negative is harmful for AUC/F1, positive is harmful for ECE). Aggregation is dataset-unweighted mean ($n=12$ datasets for C1/C2/C4/C2+C4; $n=10$ categorical-eligible datasets for C3).}
\label{fig:crosssummary}
\end{figure*}

\begin{figure*}[!t]
\centering
\includegraphics[width=0.92\textwidth]{../results/plots/paper__calibration__delta_vs_uncal.png}
\caption{Calibration gains under severe corruption, measured as change versus uncalibrated baseline. Negative values indicate improvement for ECE and log-loss. Single severe pools C2/C4 on all datasets and C3 on categorical-eligible datasets only (\texttt{n\_cat\_cols > 0}).}
\label{fig:calsummary}
\end{figure*}

\begin{figure*}[!t]
\centering
\includegraphics[width=0.92\textwidth]{../results/plots/paper__policy__summary.png}
\caption{Policy-governance gains under severe corruption for uncalibrated scores. Bars report F1 improvement and expected-cost reduction versus fixed threshold $t=0.5$. Single severe pools C2/C4 on all datasets and C3 on categorical-eligible datasets only; uncertainty intervals are reported in Table~\ref{tab:policytable}.}
\label{fig:policysummary}
\end{figure*}

\section{Discussion}
The dominant role of structured missingness indicates that coordinated feature deletion can be more harmful than moderate numerical inconsistency in operational tabular pipelines. Compound C2+C4 yields larger degradation than either single family in most settings.

From a systems perspective, results support a three-layer governance model: (1) data-pipeline controls (completeness, domain consistency, numeric range checks), (2) calibration surveillance, and (3) threshold policy governance tied to explicit error costs. Static deployment thresholds are consistently suboptimal under corruption.

Discrimination metrics alone can understate operational risk. Several settings show moderate AUC shifts but substantial ECE and expected-cost shifts, motivating corruption-aware monitoring for decision support and expert systems.

\subsection{Practical Implications}
The empirical findings can be translated into a compact operational blueprint. First, data-quality controls should monitor completeness, category novelty, and numerical consistency before scores are consumed downstream. Second, calibration should be monitored as an ongoing reliability control rather than a one-time clean-data optimization. Third, threshold policy should be reviewed explicitly against current cost assumptions instead of being fixed permanently at deployment time.

In this benchmark, C2-like failure modes are the highest-priority target for first-line controls. That implies stronger join-key validation, null-propagation audits, and mandatory completeness checks on critical feature groups. C3 and C4 remain relevant, but C2 is the most consistent degradation source and therefore offers the largest expected return for preventive engineering effort.

\subsection{Threats to Validity}
This benchmark models structured operational corruptions but cannot cover all real-world failure modes such as concept drift, label shift, delayed feedback, or adversarial manipulation. Only binary classification tasks were included, and all corruption operators were synthetic even when operationally motivated. The calibration and threshold analyses are post hoc and do not include online adaptation or human-in-the-loop intervention studies.

Additional limitations are methodological. Although 20 paired seeds improve setting-level stability, dataset-level confirmatory scope is still bounded by 12 datasets. Model hyperparameters were fixed for comparability, which improves controlled benchmarking but may understate best-case robustness achievable through aggressive per-dataset tuning. Cost assumptions are intentionally simple in this manuscript ($C_{FP}=1$, $C_{FN}=5$); real deployments may require richer utility models with intervention budgets, delayed outcomes, and conditional costs.

\subsection{Threats to Validity}
This benchmark models structured operational corruptions but cannot cover all real-world failure modes (e.g., concept drift, label shift, delayed feedback, adversarial manipulation). Only binary classification tasks were included, and all corruption operators were synthetic even when operationally motivated. The calibration and threshold analyses are post hoc and do not include online adaptation or human-in-the-loop intervention studies. Future work should extend to longitudinal monitoring and domain-specific cost models.

Additional limitations are methodological. First, although 20 paired seeds improve setting-level stability, dataset-level confirmatory scope is still bounded by 12 datasets; larger benchmark coverage would improve external generalization of significance profiles. Second, model hyperparameters were intentionally fixed for comparability, which improves controlled benchmarking but may understate best-case robustness achievable via aggressive per-dataset tuning. Third, cost assumptions are intentionally simple in this manuscript ($C_{FP}=1$, $C_{FN}=5$); real deployments may require richer utility models with intervention budgets, delayed outcomes, and conditional costs.

Threats to external validity should be explicit. The corruption families are designed for operational plausibility, not incident-level realism for any single production system. Accordingly, transfer to a specific deployment depends on whether local fault geometry resembles C1--C4 stress channels and on whether local governance can apply comparable calibration and threshold controls. A second external-validity caveat is C1: manuscript-default C1 uses label-agnostic anchoring, while label-informed anchoring is retained only as stress-synthesis sensitivity analysis and is not used for confirmatory ranking evidence. A third caveat concerns data provenance granularity: OpenML IDs/versions are pinned, but external artifact-level content hashes are not available for OpenML-hosted records in this workflow.

These limitations do not invalidate the central reliability pattern, but they delimit scope. The results should be interpreted as a structured stress benchmark under explicit assumptions, not as a universal guarantee for all tabular deployments.

\section{Conclusions}
This study provides a reproducible stress benchmark for tabular ML reliability under structured operational corruptions. In this benchmark run, descriptive cross-dataset severe-effect averages indicate that C2 has the largest mean harmful single-family effects among non-label-informed channels, and C2+C4 has the largest mean harmful compound effects for discrimination and F1. Cluster-aware sensitivity preserves the headline harmful directions for the predefined primary effects. Calibration improves probability quality under corruption, and threshold governance yields meaningful expected-cost reductions.

Inferential outputs are strengthened by 20 paired seeds per dataset and setting, while cross-dataset inference still depends on 12 datasets. Practical interpretation therefore prioritizes effect magnitudes, confidence intervals, and explicit predefined-hypothesis p/q reporting over binary significance labels. C1 findings are retained as descriptive stress diagnostics under label-agnostic default anchoring, with side-by-side label-informed sensitivity reported to bound interpretation.

The main deployment implication is operational: robust decision support requires integrated governance of data quality, calibration, and threshold policy. Clean-test discrimination alone is not sufficient for reliability claims in production tabular systems.

\section*{Acknowledgment}
The author thanks the maintainers of OpenML and the open-source scientific Python ecosystem.

\section*{Declarations}
Funding: This research received no external funding. Conflicts of interest: The author declares no conflict of interest. Code and reproducibility materials: the scripts, generated tables, and supporting artifacts used to prepare this submission are included in the accompanying submission package.

\section*{Declaration of Generative AI in Scholarly Writing}
During the preparation of this work, the author used generative AI tools (OpenAI ChatGPT/Codex) to improve grammar and readability and to assist with code-level debugging and reproducibility checks. All generated suggestions were reviewed, tested, and edited by the author, who accepts full responsibility for the final manuscript and supporting artifacts.

\bibliographystyle{unsrt}
\bibliography{references}

\end{document}

